\documentclass[10pt,a4paper]{amsart}
\usepackage[margin=19mm]{geometry}
\usepackage{amsmath,amssymb,mathtools}
\usepackage[scr=rsfs,cal=euler]{mathalfa}
\usepackage{xfrac}
\usepackage[unicode,hidelinks,bookmarks=false]{hyperref}
\newcommand{\ie}{i.e.\ }
\newcommand{\cf}{cf.\ }
\newcommand{\resp}{respectively\ }
\newcommand{\Subsection}{\subsection}
\providecommand{\nobreakdash}{\nobreak\hbox{-}}
\setlength{\emergencystretch}{2em}
\multlinegap0pt
\newcommand{\R}{\mathbb R}
\DeclareMathOperator{\clust}{clust}
\DeclareMathOperator{\dist}{dist}
\DeclareMathOperator{\dom}{dom}
\theoremstyle{plain}
\newtheorem{theorem}{Theorem}[section]
\newtheorem{lemma}[theorem]{Lemma}
\newtheorem{proposition}[theorem]{Proposition}
\newtheorem{corollary}[theorem]{Corollary}
\theoremstyle{definition}
\newtheorem{assumption}[theorem]{Assumption}
\newtheorem{definition}[theorem]{Definition}
\newtheorem{algorithm}[theorem]{Algorithm}
\theoremstyle{remark}
\newtheorem{remark}[theorem]{Remark}
\newtheorem{example}[theorem]{Example}
\title[Whole-Sequence Convergence of Variable-Smoothing Full-Splitting Methods via a Lifted Kurdyka-\L{}ojasiewicz Framework]{Whole-Sequence Convergence of Variable-Smoothing Full-Splitting Methods via a Lifted Kurdyka-\L{}ojasiewicz Framework}
\author{Min Tao}
\date{Source: arXiv:2609.28306v1 (CC BY 4.0)}
\begin{document}
\maketitle

\noindent\emph{Source and licence.} The delivered work is the article shown in the title above, version arXiv:2609.28306v1 (CC BY 4.0), distributed under the Creative Commons Attribution 4.0 International licence, as recorded in the arXiv licence metadata for this exact version and shown on the abstract page saved with this item. The full licence notice, which carries the licence URL and the address of that abstract page, is the bundled file \texttt{licenses/arxiv-2609.28306-CC-BY-4.0.txt}.
\par\smallskip

\noindent\textit{Editorial scope.} Counted unit: the article's Theorem (source label lem:scaled-KL) (source0.tex lines 1585--1606, source label \texttt{lem:scaled-KL}): ``Let be a proper lower semicontinuous function, let be bounded, and write . Suppose that and that satisfies on . Assume that has the KL property on . Let , and let be a nondecreasing positive scalar se''. It is delivered together with the article's complete original proof (source0.tex lines 1607--1665, one proof environment adjacent to the statement, 59 lines), copied line for line from the article's own text. Required context retained uncounted: eq:KL-def (lines 694--696); lem:uniform-KL (lines 701--709). Editorial formatting only: an AMS \texttt{amsart} preamble that restates the article's own macro definitions and its own theorem declarations, and the theorem environments the retained lines use; the article's own class file is neither shipped nor input; the retained environments are renumbered sequentially in order of appearance, whereas the article prints them with its own numbering; the article's equation numbering is not reproduced and every cross-reference inside the retained text resolves to the wrapper's numbering. No citation occurs inside any retained line, so the wrapper carries no bibliography; All retained bodies are exact copies of the bundled \texttt{sources/211574/source0.tex}, so no omission receipt applies. No asset-inclusion command occurs inside any retained span and the package ships no image asset.


\section*{Retained context: eq:KL-def (source0.tex lines 694--696)}

\begin{eqnarray}\label{eq:KL-def}
 \varphi'(f(w)-f(\overline w))\dist(0,\partial f(w))\geq1
\end{eqnarray}


\section*{Retained context: lem:uniform-KL (source0.tex lines 701--709)}

\begin{lemma}[Uniformized KL inequality]\label{lem:uniform-KL}
Let $f$ be proper and lower semicontinuous and let
$\Omega\subseteq\dom\partial f$ be compact. If $f$ is constant on
$\Omega$ and has the KL property at every point of $\Omega$, then
there are $\varepsilon,a>0$ and a concave desingularizing function
$\varphi$ such that~\eqref{eq:KL-def}, with $f(\overline w)=f(\Omega)$,
holds whenever $\dist(w,\Omega)<\varepsilon$ and
$f(\Omega)<f(w)<f(\Omega)+a$.
\end{lemma}


\section{The counted unit: Theorem (source label lem:scaled-KL) and its complete original proof}

\begin{theorem}\label{lem:scaled-KL}
Let $H:\R^d\to(-\infty,+\infty]$ be a proper lower semicontinuous
function, let $\{U^k\}_{k\ge 1}\subset\R^d$ be bounded, and write
$H_k:=H(U^k)$. Suppose that $H_k\to H_\infty\in\R$ and that
$\Omega:=\clust\{U^k\}\subseteq\dom\partial H$ satisfies
$H\equiv H_\infty$ on $\Omega$. Assume that $H$ has the KL property
on $\Omega$. Let $s_k\geq0$, and let $\{\delta_k\}$ be a
nondecreasing positive scalar sequence. Suppose that, for all
sufficiently large $k$,
\begin{eqnarray*}
\begin{aligned}
\textnormal{(a)}\quad&H_k-H_{k+1}\geq a\delta_ks_k^2,\\
\textnormal{(b)}\quad&\dist(0,\partial H(U^{k+1}))
 \leq b(\delta_ks_k+\varepsilon_k),
\end{aligned}
\end{eqnarray*}
where $a,b>0$ and $\varepsilon_k\geq0$, with
\begin{eqnarray}\label{eq:abstract-error-sum}
\sum_{k=1}^{+\infty}\frac{\varepsilon_k}{\delta_{k+1}}<\infty.
\end{eqnarray}
Then $\sum_{k=1}^{+\infty} s_k<\infty$.
\end{theorem}

\begin{proof}
A bounded sequence approaches its compact cluster set. Thus, Lemma \ref{lem:uniform-KL} provides a  concave desingularizing function \(\varphi\) such that, for all sufficiently large \(k\) satisfying
$H_{k+1}>H_\infty$,
\begin{eqnarray}\label{eq:uniform-KL}
 \varphi'(H_{k+1}-H_\infty)
 \dist\!\left(0,\partial H(U^{k+1})\right)\geq1.
\end{eqnarray}

Suppose first that $H_j=H_\infty$ for some sufficiently large $j$.
Monotonicity and convergence then give
$H_k=H_\infty$ for every $k\geq j$, and the descent estimate implies
$s_k=0$ for every $k\geq j$. The conclusion follows in this case.

It remains to consider the case in which $H_k>H_\infty$ when $k$ is sufficiently large. Define
\[
 \Delta_{k+1}:=
 \varphi(H_{k+1}-H_\infty)
 -\varphi(H_{k+2}-H_\infty)\geq0.
\]
Concavity of $\varphi$, the descent estimate at index $k+1$,
\eqref{eq:uniform-KL}, and the relative-error estimate at index $k$
yield
\begin{eqnarray}\label{eq:Delta-lower}
\begin{aligned}
 \Delta_{k+1}
 &\geq\varphi'(H_{k+1}-H_\infty)(H_{k+1}-H_{k+2})\\
 &\geq\frac{a\delta_{k+1}s_{k+1}^2}
 {b(\delta_ks_k+\varepsilon_k)}.
\end{aligned}
\end{eqnarray}
The denominator in~\eqref{eq:Delta-lower} is positive. Indeed, if
$\delta_ks_k+\varepsilon_k=0$, then the relative-error estimate gives
$\dist(0,\partial H(U^{k+1}))=0$, contradicting
\eqref{eq:uniform-KL}.

Set $C:=b/a$. Since $\delta_k\leq\delta_{k+1}$,
$$
 s_{k+1}^2
 \leq C\left(
 \frac{\delta_k}{\delta_{k+1}}s_k
 +\frac{\varepsilon_k}{\delta_{k+1}}\right)\Delta_{k+1}
 \leq C\left(s_k+\frac{\varepsilon_k}{\delta_{k+1}}\right)
 \Delta_{k+1}.$$
The inequality $2\sqrt{uv}\leq u+v$ therefore gives
\begin{eqnarray}\label{eq:length-recursion}
 2s_{k+1}\leq s_k+\frac{\varepsilon_k}{\delta_{k+1}}
 +C\Delta_{k+1}.
\end{eqnarray}
Summing~\eqref{eq:length-recursion} from a sufficiently large index
$N$ to $T\geq N$ and telescoping the terms $\Delta_{k+1}$ yield
\begin{eqnarray*}
 \sum_{k=N}^{T}s_{k+1}+s_{T+1}
 \leq s_N+\sum_{k=N}^{T}\frac{\varepsilon_k}{\delta_{k+1}}
 +C\varphi(H_{N+1}-H_\infty).
\end{eqnarray*}
Letting $T\to\infty$ and using~\eqref{eq:abstract-error-sum}, we obtain
$\sum_{k=N}^{\infty}s_{k+1}<\infty$. Adding the finite initial
segment completes the proof.
\end{proof}

\end{document}
